\chapter{Clostridioides difficile Testing}

\section*{What Is This Test?}
\textit{Clostridioides difficile} testing detects toxigenic \textit{C. difficile}, an anaerobic, spore-forming, gram-positive bacillus that is the leading cause of healthcare-associated infectious diarrhea and antibiotic-associated colitis. \textit{C. difficile} infection (CDI) occurs when the normal gut microbiota is disrupted by antibiotics (especially clindamycin, fluoroquinolones, cephalosporins, and broad-spectrum penicillins), allowing ingested or resident \textit{C. difficile} spores to germinate, colonize the colon, and produce toxins. The major virulence factors are toxin A (TcdA, an enterotoxin) and toxin B (TcdB, a cytotoxin), encoded by the \textit{tcdA} and \textit{tcdB} genes within the pathogenicity locus (PaLoc). Some hypervirulent strains (ribotype 027/NAP1/BI) produce a third toxin (binary toxin/CDT) and are associated with higher toxin production, increased sporulation, more severe disease, and higher mortality. CDI clinical spectrum ranges from mild, self-limited diarrhea to severe pseudomembranous colitis to fulminant colitis with toxic megacolon, ileus, colonic perforation, and sepsis. CDI recurs in 15--25\% of patients after a first episode and in 40--65\% after two or more episodes. The diagnostic approach involves a combination of tests, with the most widely accepted algorithm being a 2-step strategy: initial screening with a highly sensitive test (NAAT PCR or GDH antigen immunoassay), followed by a toxin immunoassay to confirm the presence of free toxin in stool.

\section*{Alternative Names}
\begin{itemize}
  \item C. diff Test
  \item C. difficile Toxin
  \item CDI Testing
  \item \textit{Clostridioides difficile} PCR
  \item Stool C. difficile
  \item C. diff GDH + Toxin EIA
  \item C. difficile Cytotoxin Assay
  \item CDI NAAT
\end{itemize}
\section*{Principle}
Nucleic acid amplification test (NAAT/PCR): Real-time PCR targeting the \textit{tcdB} gene (or \textit{tcdA}) of toxigenic \textit{C. difficile}. Platforms include Cepheid GeneXpert C. difficile (results in 45 minutes), Roche cobas Cdiff, BD MAX Cdiff, and BioFire FilmArray GI Panel (multiplex PCR). Sensitivity 93--97\%, specificity 95--98\%. PCR detects the gene encoding toxin production, but does NOT confirm that toxin protein is being produced in the stool. This is a key limitation: PCR is highly sensitive but may detect patients colonized with toxigenic \textit{C. difficile} who do NOT have active CDI but are carriers (especially in hospitalized patients, where colonization rates are 10--20\%). NAAT alone (without toxin testing) leads to overdiagnosis of CDI. Glutamate dehydrogenase (GDH) antigen immunoassay: EIA detecting GDH, a cell wall-associated enzyme produced in large quantities by all \textit{C. difficile} strains (both toxigenic and nontoxigenic). Highly sensitive (95--99\%) but nonspecific for toxin production; cannot distinguish toxigenic from nontoxigenic \textit{C. difficile}. A negative GDH has a very high NPV (>99\%) for ruling out CDI. Toxin A/B enzyme immunoassay (EIA): Detects TcdA and TcdB toxin proteins in stool by ELISA. Specific for active toxin production but has low sensitivity (50--70\%) compared to NAAT and cytotoxin assay. A positive toxin EIA confirms active CDI with high specificity. Cell culture cytotoxin neutralization assay (CCNA): Stool filtrate is inoculated onto fibroblast monolayers; if \textit{C. difficile} toxin is present, it causes cell rounding (cytopathic effect) that is neutralized by antitoxin. The historical gold standard. Sensitivity 67--80\%, specificity >99\%. Labor-intensive, slow (24--48 hours), and now rarely performed outside reference laboratories.

The recommended 2-step algorithm: (1) Screening with GDH EIA and/or NAAT (highly sensitive); if negative, CDI is essentially excluded. (2) If screening is positive (GDH antigen detected and/or \textit{tcdB} PCR positive), reflex to a toxin A/B EIA. GDH+/toxin+ OR PCR+/toxin+ confirms CDI. GDH+/toxin- OR PCR+/toxin- is ambiguous: may represent \textit{C. difficile} colonization without active CDI (carrier state) or CDI with toxin levels below the EIA detection limit. In these patients, clinical assessment is critical (are they having clinically significant diarrhea? Are there alternative explanations?). The CDC recommends against testing formed stool (Bristol stool types 1--4), as asymptomatic colonization is common and a positive result in the absence of diarrhea does not indicate CDI.

\section*{History}
\textit{Clostridioides difficile} was first isolated from the stool of healthy neonates in 1935 by Hall and O'Toole (hence ``difficile,'' difficult to culture) \cite{hall1935intestinal}. The association between antibiotic-associated pseudomembranous colitis and \textit{C. difficile} was established in the late 1970s by John Bartlett and colleagues. The \textit{C. difficile} cytotoxin was identified as the causative factor in 1978. Metronidazole and vancomycin were established as effective treatments in the 1980s. The BI/NAP1/027 epidemic strain emerged in the early 2000s in North America and Europe, causing severe outbreaks with increased morbidity and mortality. PCR for detection of \textit{tcdB} became commercially available in the 2000s and was rapidly adopted. However, by the 2010s, recognition grew that PCR alone (without toxin testing) led to overdiagnosis because of the high rate of asymptomatic colonization detected by sensitive PCR. In response, IDSA/SHEA 2017 CDI guidelines recommended a 2-step algorithm (GDH + toxin or NAAT + toxin) rather than NAAT alone \cite{mcdonald2018clinical}. Fidaxomicin (2011) and bezlotoxumab (2016, a monoclonal antibody against toxin B) expanded therapeutic options. Fecal microbiota transplantation (FMT) became an accepted treatment for multiply recurrent CDI, with cure rates of 80--90\%.

\section*{Why This Test Is Ordered}
\begin{itemize}
  \item Diagnosis of \textit{C. difficile} infection in patients with new-onset, clinically significant diarrhea ($\geq$3 unformed/liquid stools in 24 hours) who are hospitalized, recently hospitalized (within 12 weeks), on antibiotics (within 12 weeks), or have risk factors (age $\geq$65, immunocompromised, PPI use---though the PPI association is debated)
  \item Investigation of healthcare-associated or antibiotic-associated diarrhea when other causes (enteral feeding, medications, other infections) have been excluded or are unlikely
  \item Evaluation for CDI in immunocompromised patients (HSCT, SOT, chemotherapy, HIV) with otherwise unexplained diarrhea, as CDI is more common and more severe in these populations
  \item CDI should be suspected in patients on antibiotics or within 8--12 weeks of antibiotic exposure who develop abdominal pain, leukocytosis (WBC often 15,000--50,000+), ileus, or toxic megacolon (CDI may present without diarrhea in the setting of ileus!)
  \item Follow-up testing is NOT recommended to document cure: toxin detection can persist for weeks after successful treatment in cured but colonized patients. Test of cure leads to false-positive diagnoses and unnecessary treatment
\end{itemize}

\section*{Primary vs Secondary Test}
\textit{C. difficile} testing by the 2-step algorithm (GDH/NAAT screening followed by toxin EIA) is the primary diagnostic approach. No single test is optimal alone. Stool should ONLY be tested from patients with clinically significant diarrhea (Bristol stool types 5--7). Testing formed stool from asymptomatic patients is strongly discouraged (high rate of positive results due to colonization, which should not be treated). Empiric therapy for fulminant CDI (ileus, toxic megacolon, shock) should commence before test results are available.

\section*{How This Test Is Performed}
Stool specimen: A fresh, unformed (liquid or semisolid) stool specimen is collected in a clean, dry, leak-proof container. Approximately 5--10 mL of liquid stool is sufficient. The specimen should be transported to the laboratory within 2 hours at room temperature, or refrigerated at 2--8\textdegree{}C and tested within 24--48 hours (toxin degrades over time at room temperature). Specimens for PCR are more stable than those for toxin EIA. Formed stool (Bristol types 1--4) should be REJECTED by the laboratory per IDSA/SHEA guidelines, with a comment that CDI testing should be performed only on unformed stool. One specimen is sufficient for diagnosis (sensitivity is not improved with repeat testing). If the initial stool specimen is negative but CDI remains clinically suspected, a second specimen may be tested (but sensitivity of the first specimen is >90\%). Testing more than 2 specimens is not recommended. For patients with ileus and no diarrhea, a rectal swab can be used for NAAT.

\section*{What Preparation Is Needed}
No patient preparation. Patients should not receive barium, enemas, or bowel preparations before stool collection, as these can alter the specimen. Laxatives and stool softeners should be avoided if possible, as they may confound the assessment of diarrhea. Antibiotics administered for the current episode should be noted but do not preclude testing.

\section*{Contraindications and Precautions}
\begin{itemize}
  \item The most common error in CDI diagnosis is testing patients who do NOT have clinically significant diarrhea (Bristol types 1--4). In hospitalized patients, 10--20\% are asymptomatically colonized with toxigenic \textit{C. difficile}. Testing formed stool in these patients yields positive results that lead to unnecessary antibiotic therapy, increased costs, and potential adverse effects with zero clinical benefit. TEST ONLY UNFORMED STOOL. A positive NAAT alone (PCR+/toxin-) is insufficient to diagnose CDI definitively; clinical correlation and toxin EIA results are essential. NAAT+/toxin- patients with clinically significant diarrhea and no alternative explanation may have CDI (particularly if they have been on antibiotics and have leukocytosis). Conversely, NAAT+/toxin- patients with mild or atypical symptoms may be carriers with another cause of diarrhea; clinical judgment is paramount. Test of cure is NOT recommended because
  \item (1) toxin EIA can remain positive for weeks after successful treatment;
  \item (2) PCR can remain positive for >30 days; and
  \item (3) clinical resolution, not microbiologic eradication, is the treatment goal. Retesting should be limited to patients with recurrent symptoms $\geq$14 days after treatment completion.
\end{itemize}
\section*{Risks}
No risks from stool collection. The primary risk of CDI testing is overdiagnosis (testing formed stool or testing asymptomatic carriers), which leads to unnecessary vancomycin/fidaxomicin exposure and selection for VRE and other resistant organisms. Undiagnosed CDI carries significant morbidity: progression to pseudomembranous colitis, toxic megacolon, colectomy, and death (all-cause mortality attributable to CDI is 5--15\%).

\section*{What Happens After the Test}
GDH/NAAT results available in 1--4 hours (rapid PCR in 45--60 minutes). Toxin EIA results in 1--3 hours. Cytotoxin assay results in 24--48 hours. Algorithm interpretation:
  - \textbf{GDH negative or PCR negative:} CDI is essentially excluded (NPV >99\%). Evaluate for alternative causes of diarrhea. Discontinue empiric anti-\textit{C. difficile} antibiotics if started
  - \textbf{GDH positive + toxin EIA positive:} Confirmed CDI. Initiate treatment. Discontinue the inciting antibiotic(s) if possible; this is critical for resolution
  - \textbf{GDH positive or PCR positive + toxin EIA negative:} Ambiguous. Clinical assessment: Does the patient have ongoing clinically significant diarrhea (Bristol 5--7, $\geq$3 episodes/24h)? Is the WBC elevated? Are there alternative explanations for diarrhea (enteral tube feeds, laxatives, other infections)? If high clinical suspicion for CDI, treat. If mild symptoms and an alternative cause, consider observation. The NAAT cycle threshold (Ct) value, if available, may help: low Ct values (<28 cycles) correlate with high organism burden and higher likelihood of true CDI; high Ct values correlate with colonization

\section*{Understanding Results}

\subsection*{Below Normal}
\begin{itemize}
  \item \textbf{GDH antigen negative:} No \textit{C. difficile} antigen detected. Essentially excludes \textit{C. difficile} carriage and CDI (NPV >99\%). No further testing for \textit{C. difficile} is needed
  \item \textbf{NAAT/PCR negative for \textit{tcdB}:} \textit{C. difficile} toxin B gene not detected. Same high NPV. CDI excluded
  \item \textbf{Toxin A/B EIA negative in the setting of positive GDH or PCR:} Indicates possible \textit{C. difficile} colonization without active toxin production, or active CDI with very low toxin levels below the EIA detection limit. Clinical correlation required. If the patient has ongoing clinically significant diarrhea, treating as presumptive CDI may be appropriate, especially if other causes are excluded
  \item \textbf{Negative testing in a patient with fulminant colitis/ileus:} In the setting of ileus, stool may be absent or testing may be falsely negative; empiric treatment for CDI and surgical consultation for possible toxic megacolon should NOT be delayed by negative test results
\end{itemize}

\subsection*{Above Normal}
\begin{itemize}
  \item \textbf{GDH positive AND Toxin A/B EIA positive:} Confirmed CDI. Treat with vancomycin 125 mg PO QID x 10 days, or fidaxomicin 200 mg PO BID x 10 days (preferred for patients at high risk for recurrence: age $\geq$65, immunocompromised, history of CDI, severe CDI). Metronidazole 500 mg PO TID x 10 days is an alternative for non-severe CDI when vancomycin or fidaxomicin are not available. Do not use metronidazole beyond the first recurrence (inferior efficacy compared to vancomycin)
  \item \textbf{PCR positive for \textit{tcdB} AND Toxin A/B EIA positive:} Confirmed CDI (same as GDH+/toxin+). Treat accordingly
  \item \textbf{NAAT positive (PCR+) with low Ct value:} High organism burden. Correlates more strongly with true CDI than colonization. Toxin EIA is typically positive at low Ct values
  \item \textbf{NAAT positive for NAP1/027 ribotype (GeneXpert):} The hypervirulent NAP1/BI/027 strain is associated with higher toxin production, increased recurrence rates, and higher mortality. In a patient with confirmed CDI, identifying this strain should prompt careful monitoring and consideration of fidaxomicin as first-line therapy (fidaxomicin may reduce recurrence rates compared to vancomycin). Contact precautions with soap-and-water hand hygiene (alcohol sanitizer does NOT kill \textit{C. difficile} spores) should be rigorously enforced
\end{itemize}

\section*{Normal Results}
\textbf{GDH antigen negative; NAAT negative; Toxin A/B EIA negative.} No evidence of \textit{C. difficile} infection or colonization.

\section*{Follow-up Tests}
\begin{itemize}
  \item \textbf{Repeat \textit{C. difficile} testing:} ONLY indicated for recurrent symptoms $\geq$14 days after a completed treatment course and symptom resolution. Persistent testing during or immediately after treatment is not recommended (false positives due to persistent nucleic acid or antigen shedding without active disease)
  \item \textbf{Abdominal imaging (CT abdomen/pelvis with contrast):} For patients with severe CDI (ileus, abdominal distension, severe pain, peritonitis). Findings may include colonic wall thickening, pericolonic stranding, ascites, and, in severe cases, toxic megacolon (colonic dilatation >6 cm with systemic toxicity) or pneumatosis. CT findings correlate with disease severity
  \item \textbf{Complete blood count and metabolic panel:} For assessing CDI severity. WBC >15,000/$\mu$L, creatinine >1.5x baseline, serum albumin <3 g/dL, and lactate elevated suggest severe disease. WBC >35,000 or lactate >2.2 mmol/L are particularly concerning
  \item \textbf{Flexible sigmoidoscopy or colonoscopy:} May be performed when the diagnosis is in doubt, CDI is suspected but stool testing is negative, or there is an atypical presentation. Pseudomembranes (yellow-white plaques 2--10 mm on an erythematous, edematous mucosa) are pathognomonic. Biopsy confirms pseudomembranous colitis. Avoid full colonoscopy with insufflation in severe disease (risk of perforation)
  \item \textit{\textbf{C. difficile}} \textbf{strain typing (PCR ribotyping or PFGE):} Performed at reference laboratories during outbreak investigations to confirm clonality and identify the strain (027, 078, etc.)
  \item \textbf{Fecal microbiota transplantation (FMT):} Considered for patients with $\geq$3 episodes of recurrent CDI after a pulsed-tapered vancomycin course. FMT restores gut microbiota diversity with cure rates of 80--90\%. FDA-approved live biotherapeutic products (Rebyota, Vowst) are now available as standardized alternatives to donor stool FMT
\end{itemize}

\section*{References}
\bibliographystyle{unsrtnat}
\bibliography{references}

\clearpage
\section*{Regulatory Status and Authorization Date}
This chapter describes a test category, not a specific commercial product. FDA, CDSCO, EU IVDR, and MHRA approval or registration dates therefore cannot be assigned without the assay manufacturer, product name, instrument, and intended use. For laboratory-developed tests, an individual agency approval date may not exist. See the Preface for the interpretation of regulatory status.

\clearpage
